\documentclass[10pt,a4paper]{article}

\usepackage[utf8]{inputenc}
\usepackage[T1]{fontenc}
\usepackage[english]{babel}
\usepackage[margin=2cm]{geometry}
\usepackage{graphicx}
\usepackage{hyperref}

% Placeholder box for a figure that still has to be produced by hand.
% #1 = caption; the body describes what belongs in it.
\newsavebox{\placebox}
\newenvironment{placeholder}[1]
  {\par\medskip\noindent\begin{center}
   \setlength{\fboxsep}{6pt}%
   \begin{lrbox}{\placebox}%
   \begin{minipage}{0.88\linewidth}
   \centering\textbf{[ PLACEHOLDER: #1 ]}\\[0.8em]
   \raggedright\footnotesize}
  {\end{minipage}\end{lrbox}\fbox{\usebox{\placebox}}\end{center}\medskip}

\title{\textbf{mapping\_pathfinding}\\[0.3em]
       \large Mapping, Localization and Path Planning for the Duckiebot City Challenge}
\author{Nico Sander}
\date{\today}
\usepackage{titlesec}
\titlespacing*{\section}{0pt}{1.2ex plus .5ex}{0.8ex}
\titlespacing*{\subsection}{0pt}{1ex plus .3ex}{0.5ex}
\setlength{\parskip}{0.4ex}

\begin{document}
\maketitle

\section{System Overview}

\texttt{mapping\_pathfinding} is a ROS 1 (Noetic) package that lets a Duckiebot
autonomously explore a small city of streets and intersections, build a map of
it, and afterwards drive a sequence of announced gates in the fastest order it
can find.

The mission is split into two phases:

\begin{description}
  \item[MAPPING (untimed)] The bot drives every street of the city at least
    once. While driving it records which \emph{gate} (an AprilTag mounted over
    a street) belongs to which street, and it measures how long each street and
    each turn actually takes.
  \item[GATE\_RUN (timed)] The gate order is announced on site as a list of tag
    IDs. The bot is placed at a known street, plans a route through the gates in
    that order using the map and the timings from the mapping phase, and drives
    it.
  \item[DONE] After the last gate the bot stops and reports the elapsed time.
\end{description}

The city is modelled as a graph: nodes are intersections, edges are streets.
Every intersection has up to four numbered \emph{ports} (1--4) arranged
cyclically, so that for a bot entering through port $p$, port $p+1$ is on its
right and $p-1$ on its left. Because U-turns are forbidden, planning does not
happen over nodes but over directed states \texttt{(node, entry\_port)} --
``approaching this intersection, having entered through this port''. This state
also identifies the street that was just driven, which is what allows the
planner to target individual streets and gates. The map itself lives in a single
JSON file (\texttt{config/city.json}) that every node reads, so there is only
ever one definition of the city.

The package is deliberately split into thin ROS nodes and a ROS-free core. All
graph, planning, timing and filtering logic is implemented in plain Python
modules without any ROS import, which means the whole decision-making core is
covered by offline unit tests (currently 209) that run without a robot.

\begin{placeholder}{ROS node graph}
  A block diagram of the running ROS graph. Draw one box per node
  (\texttt{detect\_lane}, \texttt{detect\_signs}, \texttt{detect\_intersection},
  \texttt{switch\_control}, \texttt{control\_wheels},
  \texttt{mapping\_pathplanning}, \texttt{visualization}, \texttt{dashboard})
  plus the camera and the wheel driver as external boxes, and connect them with
  arrows labelled by the topic name. The important edges are:
  camera $\rightarrow$ \texttt{detect\_lane} and \texttt{detect\_signs};
  \texttt{/detect/lane\_borders} $\rightarrow$ \texttt{detect\_intersection};
  \texttt{/detect/intersection} and \texttt{/detect/sign} $\rightarrow$
  \texttt{switch\_control}; \texttt{/switch/mode} and
  \texttt{/switch/turn\_direction} $\rightarrow$ \texttt{control\_wheels}
  \emph{and} $\rightarrow$ \texttt{mapping\_pathplanning};
  \texttt{/detect/sign\_detections} $\rightarrow$
  \texttt{mapping\_pathplanning}; and the edge that closes the loop,
  \texttt{/plan/turn\_command} $\rightarrow$ \texttt{switch\_control}.
  Also show \texttt{/mapping/state} $\rightarrow$ \texttt{visualization}.
  Worth marking visually: the three layers -- perception (left), behaviour
  (middle), mission/planning (right) -- and the fact that the mission layer
  feeds back into the behaviour layer. A screenshot of \texttt{rqt\_graph} is an
  acceptable substitute, but a hand-drawn version is far more readable because
  it can leave out the debug topics.
\end{placeholder}

\section{Key Components}

\subsection{Perception}

\paragraph{\texttt{detect\_lane.py}} Segments the camera image with a PyTorch
U-Net into white lane markings, the yellow centre line and red stop lines. From
the segmentation it computes a cross-track error -- how far the bot is off the
centre of its lane -- and publishes it on \texttt{/detect/lane}. It also
publishes the geometric lane borders (including the position of a red line in
the image) on \texttt{/detect/lane\_borders}, plus debug images for the
dashboard. The whole perception pipeline runs at 30\,Hz on the laptop, around
20\,ms per frame against a 33\,ms budget.

\paragraph{\texttt{detect\_signs.py}} Detects AprilTags with
\texttt{pupil\_apriltags} and publishes on two topics, because the two consumers
need different things. \texttt{/detect/sign} carries only the ID of the closest
tag and is unthresholded -- \texttt{switch\_control} needs to see intersection
signs at any size. \texttt{/detect/sign\_detections} carries \emph{every}
detection with its corners, its area in pixels and an accept/reject decision;
the mission node needs the area to tell whether a gate is on the current street
or somewhere beyond the intersection.

\paragraph{\texttt{detect\_intersection.py}} A small logic node that turns the
lane geometry into an intersection state: no intersection, intersection
approaching, or stop line reached. It does this purely from the vertical
position of the detected red line, without touching the image again.

\subsection{Behaviour and Actuation}

\paragraph{\texttt{switch\_control.py}} The behaviour state machine. It decides
\emph{what the bot is currently doing} -- lane following, approaching an
intersection, stopped at the red line, or crossing -- and publishes that as
\texttt{/switch/mode} together with the turn direction on
\texttt{/switch/turn\_direction}. The turn direction normally comes from the
mission node over \texttt{/plan/turn\_command}; if that command is stale or
missing, the node falls back to its older behaviour of picking a direction from
the intersection sign, so a dead mission node degrades gracefully instead of
stopping the robot. Crossings are closed-loop: the manoeuvre ends when lane
markings are visible again, not when a timer expires -- the configured duration
is only a timeout. This is what makes the stack tolerant of the bot not being
perfectly square at the stop line.

\paragraph{\texttt{control\_wheels.py}} The only node that talks to the motors.
It translates the current drive mode into wheel commands: in lane following it
runs a PID controller on the cross-track error (with a low-pass filtered
derivative term so it stays stable at 30\,Hz), and during a crossing it plays
back a fixed arc for the commanded direction until \texttt{switch\_control} ends
the manoeuvre. Measured latency from photon to command is roughly 75\,ms.

\subsection{Mission, Mapping and Planning}

\paragraph{\texttt{mapping\_pathplanning.py}} The node that owns the mission and
the largest single component of the package. It does four things. \emph{(1)
Localization:} graph dead-reckoning -- the believed position advances by one
street every time \texttt{switch\_control} reports a completed crossing.
\emph{(2) Mapping:} it writes gates onto the streets they were seen on. \emph{(3)
Planning:} after every crossing it replans the whole remaining route, which is
cheap on a city this size and means a corrected position immediately produces a
corrected route. \emph{(4) Commanding:} it publishes the next turn on
\texttt{/plan/turn\_command}, which closes the control loop. It also publishes
the full mission state (map, gates, path, position, phase) as JSON on
\texttt{/mapping/state} and offers a \texttt{start\_gate\_run} service for the
phase switch.

\paragraph{\texttt{planner.py} (ROS-free)} Dijkstra over the directed states
\texttt{(node, entry\_port)}. U-turns are not forbidden by a check but are
structurally impossible to express, because no direction table maps a port onto
itself. Edges are weighted with an estimated duration -- a move costs
\texttt{stop + turn + street} -- so the search prefers genuinely faster routes,
not just shorter ones; a right turn is markedly cheaper than a left one. For the
mapping phase the same module drives a greedy nearest-unvisited-street
exploration.

\paragraph{\texttt{city\_map.py} (ROS-free)} The graph itself: port geometry,
loading and validating \texttt{city.json}, and \texttt{GraphMap}, the live map
state that tracks which streets have been driven and which gate sits on which
street. Edge keys are direction-independent (\texttt{A1\_\_B1} from either end),
so ``a gate spans the whole street'' falls out for free.

\paragraph{\texttt{gate\_detection.py} (ROS-free)} Decides which gate sightings
may be believed. A gate is written onto a street permanently -- the first write
wins -- because the camera sees past the intersection and later sightings are
taken from further away and at a worse angle. Three filters guard that
irreversible write: mapping is disabled while stopped and while crossing; the
tag's area must exceed a threshold calibrated from recorded runs
(\texttt{gate\_min\_area}, currently 900\,px$^2$); and the same tag must be seen
on the same street for several consecutive frames.

\paragraph{\texttt{run\_timing.py} (ROS-free)} Measures how long each street and
each turn actually takes, using the drive-mode transitions published by
\texttt{switch\_control} as the single time base, so the measured pieces tile a
move with no gap or overlap. The mapping phase is untimed, so measuring there is
free; the timed gate run is then planned against real numbers instead of a flat
guess. Repeated measurements of a street are reduced with the median.

\paragraph{\texttt{crossing.py} and \texttt{graph\_layout.py} (ROS-free)} Two
small support modules: the first decides when an intersection crossing is over
(lane markings visible again, after a minimum blind time), the second computes
the drawing geometry for the map, with the tested property that drawn streets
must never cross each other outside a node.

\subsection{Operation and Visualization}

\paragraph{\texttt{visualization.py}} A live window showing the three things the
challenge asks to be made visible: the map that was built (streets with the
gates found on them), the planned path (dashed overlay with step numbers), and
where the bot believes it is (current street highlighted, ring on the next
intersection).

\paragraph{\texttt{dashboard.py}} A perception debug window that tiles the raw
camera image, the segmentation masks and the AprilTag bounding boxes with their
measured areas into a single view.

\paragraph{\texttt{mission\_tui.py} / \texttt{mission\_setup.py}} A menu-driven
front end for the whole session. Its structural point is that \emph{one}
\texttt{roslaunch} stays alive across both phases -- the map, the gates and the
measured times live in the mission node, and restarting it between phases would
throw them away. The TUI therefore owns the launch process and switches phases
through parameters and the service call. \texttt{mission\_setup.py} validates
the typed answers (above all the start placement, which is checked against the
city graph before anything is launched) and builds the launch command.

\section{State Diagrams and Flowcharts}

\begin{placeholder}{Mission phase state diagram}
  The top-level mission state machine of \texttt{mapping\_pathplanning.py}, with
  three states: \textbf{MAPPING}, \textbf{GATE\_RUN} and \textbf{DONE}.
  Transitions to draw: MAPPING $\rightarrow$ GATE\_RUN, labelled ``all streets
  driven \emph{and} all gates seen $\rightarrow$ halt at the red line, wait for
  the \texttt{start\_gate\_run} service''; GATE\_RUN $\rightarrow$ DONE,
  labelled ``last gate in the announced order passed''. Mark clearly that
  between MAPPING and GATE\_RUN there is a \emph{hard stop} during which the bot
  is physically repositioned, and that the transition out of MAPPING requires
  gates to be \emph{seen}, not just streets to be entered.
\end{placeholder}

\begin{placeholder}{Behaviour state machine (\texttt{switch\_control})}
  The drive-mode FSM: \textbf{LANE\_FOLLOWING} $\rightarrow$
  \textbf{APPROACHING} $\rightarrow$ \textbf{STOPPED} $\rightarrow$
  \textbf{CROSSING} $\rightarrow$ back to LANE\_FOLLOWING. Label each transition
  with its trigger: red line detected above the threshold; stop line reached;
  stop duration elapsed \emph{and} a turn command available; and -- the
  important one -- crossing ends when \emph{lane markings are visible again},
  with the configured turn duration acting only as a timeout. Add the
  \texttt{HALT} command as a transition into a holding state, and
  \texttt{/plan/resume\_driving} as the edge that puts a repositioned bot back
  into LANE\_FOLLOWING without performing a crossing. It is worth annotating the
  same diagram with which time interval \texttt{run\_timing.py} measures as
  ``turn time'' and which as ``street time''.
\end{placeholder}

\begin{placeholder}{Gate sighting decision flowchart}
  A flowchart for a single incoming detection from
  \texttt{/detect/sign\_detections}, ending in either ``write gate onto street''
  or ``discard''. Decision boxes in order: is the bot driving (not stopped, not
  crossing)? is the tag ID a gate at all? is its area $\geq$
  \texttt{gate\_min\_area}? has the same (street, gate) pair been seen for
  \texttt{gate\_confirm\_frames} consecutive detections? does this street
  already carry a gate (if yes: discard -- the first write is permanent)? This
  diagram is the best place to show \emph{why} each filter exists; a one-line
  annotation per box is enough.
\end{placeholder}

\begin{placeholder}{Control loop / per-intersection sequence}
  A sequence diagram over one intersection, with the lanes
  camera $\rightarrow$ \texttt{detect\_*} $\rightarrow$ \texttt{switch\_control}
  $\rightarrow$ \texttt{control\_wheels}, and \texttt{mapping\_pathplanning} as
  a parallel lane. Show the order: red line detected, stop, mission node
  replans and publishes the turn command, \texttt{switch\_control} enters
  CROSSING with that direction, \texttt{control\_wheels} drives the arc, lane
  markings reappear, \texttt{switch\_control} reports the completed crossing,
  and the mission node advances its believed position by one street and replans.
  This is the diagram that makes the closed loop -- and the 30\,Hz inner loop
  inside the outer per-intersection loop -- legible.
\end{placeholder}

\section{Configuration and Testing}

All tuned constants live in \texttt{config/config.json} (timings, PID gains,
thresholds, the planner's cost model); the map is \texttt{config/city.json} and
the cosmetic gate ID $\rightarrow$ colour mapping is \texttt{config/gates.json}.
Launch arguments override selected values for a single run.

Because the decision-making core is ROS-free, it is verified offline: port
geometry and the no-U-turn invariant, city validation, gate-sighting filters,
planner reachability from every start state, all gate orderings from all start
states, greedy coverage, full mission simulations and synthetic grid cities.
In addition, \texttt{launch/simulate.launch} runs the real mission node against
a fake driver, exercising the actual ROS layer -- topics, timers, the phase
switch -- without a camera, motors or a robot.

\end{document}
